% ===========================================================================
%  Chapitre 12 — Nombres complexes
%  PE 2026, composante GÉOMÉTRIE
% ===========================================================================
\bmtchapitre{Nombres complexes}
  {Utiliser les nombres complexes pour résoudre un problème dans le plan :
   formes algébrique, trigonométrique, exponentielle et polaire ; formules de
   Moivre et d'Euler ; équations et racines $n$-ièmes.}
  {Géométrie \textbullet\ environ 12 heures}

L'équation $x^{2} = -1$ n'a pas de solution réelle. Pendant des siècles, on
s'est contenté de le constater ; puis les algébristes italiens du seizième
siècle ont remarqué qu'en manipulant une racine carrée de $-1$ comme si elle
existait, on obtenait des résultats justes sur des équations bien réelles. Ce
chapitre officialise ce geste : on ajoute aux réels un nombre $i$ tel que
$i^{2} = -1$, et l'on vérifie que tout continue de fonctionner.

Ce qui suit dépasse largement l'algèbre. Chaque nombre complexe est un point
du plan, chaque produit est une rotation combinée à un agrandissement. Les
nombres complexes deviennent alors le langage le plus efficace pour décrire
les transformations du plan — et c'est à ce titre qu'ils reviennent au
baccalauréat, dans le problème de géométrie, chaque année sans exception.

% ---------------------------------------------------------------------------
\begin{revision}
Trigonométrie et géométrie plane de première.

\begin{enumerate}[label=\textbf{\arabic*.}, leftmargin=*, itemsep=0.35em]
  \item Donner les valeurs de $\cos\frac{\pi}{3}$, $\sin\frac{\pi}{3}$,
        $\cos\frac{\pi}{4}$, $\sin\frac{\pi}{4}$, $\cos\frac{\pi}{6}$,
        $\sin\frac{\pi}{6}$.
  \item Résoudre $x^{2}-4x+5 = 0$ dans $\R$. Que vaut le discriminant ?
  \item Développer $(a+b)^{2}$ et $(a+b)^{3}$.
  \item Rappeler les formules d'addition
        $\cos(a+b)$ et $\sin(a+b)$.
  \item Dans un repère orthonormé, calculer la distance $AB$ pour
        $A(1\,;2)$ et $B(4\,;6)$.
\end{enumerate}
\end{revision}

\begin{automatismes}
Sans calculatrice, sans brouillon.

\begin{enumerate}[label=\textbf{\alph*)}, leftmargin=*, itemsep=0.25em]
  \item $i^{2}$, $i^{3}$, $i^{4}$
  \item $(2+3i)+(1-i)$
  \item $(1+i)^{2}$
  \item $(1+i)(1-i)$
  \item Conjugué de $3-5i$
  \item $\abs{3+4i}$
  \item $\dfrac{1}{i}$
  \item $\cos\frac{\pi}{3}+i\sin\frac{\pi}{3}$ sous forme algébrique
\end{enumerate}
\end{automatismes}

\begin{corrige}[automatismes]
a) $-1$, $-i$, $1$ \quad b) $3+2i$ \quad c) $2i$ \quad d) $2$ \quad
e) $3+5i$ \quad f) $5$ \quad g) $-i$ \quad
h) $\frac12+i\frac{\sqrt3}{2}$.
\end{corrige}

% ===========================================================================
\section{Définition et forme algébrique}

\begin{definition}[ensemble des nombres complexes]
Il existe un ensemble $\C$ contenant $\R$, muni d'une addition et d'une
multiplication prolongeant celles de $\R$, et contenant un élément $i$
vérifiant $i^{2} = -1$. Tout élément de $\C$ s'écrit de façon unique
\[
  z = x+iy, \qquad x \in \R, \ y \in \R,
\]
écriture appelée \textbf{forme algébrique}. Le réel $x$ est la \textbf{partie
réelle} $\Re(z)$, le réel $y$ la \textbf{partie imaginaire} $\Im(z)$.
\end{definition}

\begin{avertissement}
La partie imaginaire est un \emph{nombre réel} : celle de $3-5i$ vaut $-5$, et
non $-5i$. Cette confusion coûte des points chaque année.
\end{avertissement}

\begin{definition}[conjugué]
Le \textbf{conjugué} de $z = x+iy$ est $\overline z = x-iy$.
\end{definition}

\begin{propriete}[règles sur le conjugué]
\[
  \overline{z+z'} = \overline z+\overline{z'}, \quad
  \overline{zz'} = \overline z\,\overline{z'}, \quad
  \overline{\left(\frac{z}{z'}\right)} = \frac{\overline z}{\overline{z'}},
  \quad z\overline z = x^{2}+y^{2}, \quad
  z+\overline z = 2\Re(z).
\]
De plus $z$ est réel si et seulement si $\overline z = z$, et imaginaire pur
si et seulement si $\overline z = -z$.
\end{propriete}

\begin{methode}{écrire un quotient sous forme algébrique}
On multiplie numérateur et dénominateur par le conjugué du dénominateur : le
nouveau dénominateur devient réel, et il ne reste qu'à séparer les deux
parties.
\end{methode}

\begin{exemple}[un quotient]
\[
  \frac{2+3i}{1-i}
  = \frac{(2+3i)(1+i)}{(1-i)(1+i)}
  = \frac{2+2i+3i+3i^{2}}{1+1}
  = \frac{-1+5i}{2}
  = -\frac12+\frac52\,i .
\]
\end{exemple}

\begin{entrainement}[forme algébrique]
\exo[\niveau{1}] Écrire sous forme algébrique $(3-i)(2+5i)$ et
$\dfrac{1+i}{2-i}$.

\exo[\niveau{1}] Déterminer $\Re(z)$ et $\Im(z)$ pour $z = (2+i)^{2}$.

\exo[\niveau{2}] Déterminer les réels $x$ et $y$ tels que
$(x+iy)(1+2i) = 5$.

\exo[\niveau{2}] Montrer que $\dfrac{z-1}{z+1}$ est imaginaire pur si et
seulement si $\abs z = 1$ et $z \neq -1$.

\exo[\niveau{3}] Calculer $i+i^{2}+i^{3}+\cdots+i^{2026}$.
\end{entrainement}

\begin{corrige}
\textbf{1.} $(3-i)(2+5i) = 6+15i-2i+5 = 11+13i$. Pour le quotient, on
multiplie par le conjugué du dénominateur :
$\dfrac{1+i}{2-i} = \dfrac{(1+i)(2+i)}{5} = \dfrac{1+3i}{5}
= \dfrac15+\dfrac35\,i$.

\textbf{2.} $(2+i)^{2} = 4+4i-1 = 3+4i$, donc $\Re(z) = 3$ et $\Im(z) = 4$.

\textbf{3.} $x+iy = \dfrac{5}{1+2i} = \dfrac{5(1-2i)}{5} = 1-2i$, d'où
$x = 1$ et $y = -2$.

\textbf{4.} Pour $z \neq -1$,
$\dfrac{z-1}{z+1} = \dfrac{(z-1)\overline{(z+1)}}{\abs{z+1}^{2}}$, et le
numérateur vaut $\abs z^{2}-1+2i\,\Im(z)$. Sa partie réelle est donc
$\abs z^{2}-1$ : le quotient est imaginaire pur exactement lorsque
$\abs z = 1$.

\textbf{5.} Les puissances de $i$ se répètent par blocs de quatre, et
$i+i^{2}+i^{3}+i^{4} = 0$. Comme $2026 = 4\times506+2$, les $2024$ premiers
termes se compensent et il reste
$i^{2025}+i^{2026} = i+i^{2} = -1+i$.
\end{corrige}

% ===========================================================================
\section{Représentation géométrique}

\begin{definition}[affixe]
Le plan étant muni d'un repère orthonormé direct $(O\,;\vecu,\vecv)$, on
associe au complexe $z = x+iy$ le point $M(x\,;y)$, appelé \textbf{image} de
$z$ ; réciproquement $z$ est l'\textbf{affixe} de $M$, notée $z_{M}$. L'affixe
du vecteur $\vect{AB}$ est $z_{B}-z_{A}$.
\end{definition}

\begin{visualisation}[le plan complexe]
\begin{center}
\begin{tikzpicture}[x=1.05cm, y=1.05cm, font=\footnotesize]
  \draw[bmtGris, -{Stealth[length=2.6mm]}] (-0.6,0) -- (4.4,0)
    node[anchor=north east] {$\vecu$};
  \draw[bmtGris, -{Stealth[length=2.6mm]}] (0,-0.6) -- (0,3.4)
    node[anchor=north west] {$\vecv$};
  \node[bmtGris, anchor=north east] at (0,0) {$O$};
  \filldraw[bmtIndigo] (3,2) circle (2pt);
  \node[bmtIndigo, anchor=south west] at (3.05,2.05) {$M(z)$};
  \draw[bmtIndigo, dashed] (3,0) -- (3,2) -- (0,2);
  \node[bmtGris, anchor=north] at (3,-0.08) {$x$};
  \node[bmtGris, anchor=east] at (-0.08,2) {$y$};
  \draw[bmtLaterite, line width=1.1pt] (0,0) -- (3,2);
  \node[bmtLaterite, anchor=south east, rotate=33.7] at (1.7,1.0) {$\abs z$};
  \draw[bmtVetiver, line width=0.9pt] (0.85,0) arc (0:33.7:0.85);
  \node[bmtVetiver, anchor=west] at (0.95,0.28) {$\theta$};
  \filldraw[bmtGris] (3,-2) circle (1.7pt);
  \node[bmtGris, anchor=north west] at (3.05,-1.95) {$\overline z$};
  \draw[bmtGris, dotted] (3,2) -- (3,-2);
\end{tikzpicture}
\end{center}

Chaque nombre complexe est un point. Sa distance à l'origine est son module,
l'angle que fait $\vect{OM}$ avec l'axe des abscisses est son argument. Le
conjugué est le symétrique par rapport à l'axe des réels : conjuguer, c'est
retourner la figure. Retenez cette double lecture — algébrique et
géométrique : tout le chapitre consiste à passer de l'une à l'autre.
\end{visualisation}

\begin{propriete}[distances et milieux]
$AB = \abs{z_{B}-z_{A}}$, et l'affixe du milieu de $[AB]$ vaut
$\dfrac{z_{A}+z_{B}}{2}$. Plus généralement, si $M$ est défini par une
relation vectorielle entre points connus, son affixe s'obtient en remplaçant
dans cette relation chaque vecteur $\vect{PQ}$ par la différence
$z_{Q}-z_{P}$.
\end{propriete}

% ===========================================================================
\section{Module, argument et forme trigonométrique}

\begin{definition}[module et argument]
Le \textbf{module} de $z = x+iy$ est $\abs z = \sqrt{x^{2}+y^{2}}$. Si
$z \neq 0$, un \textbf{argument} de $z$ est tout réel $\theta$ tel que
\[
  \cos\theta = \frac{x}{\abs z} \quad\text{et}\quad
  \sin\theta = \frac{y}{\abs z}.
\]
On écrit alors $z = \abs z\left(\cos\theta+i\sin\theta\right)$ :
c'est la \textbf{forme trigonométrique}.
\end{definition}

\begin{propriete}[règles de calcul]
\[
  \abs{zz'} = \abs z\abs{z'}, \qquad
  \arg(zz') = \arg z+\arg z' \ [2\pi],
\]
\[
  \left|\frac{z}{z'}\right| = \frac{\abs z}{\abs{z'}}, \qquad
  \arg\!\left(\frac{z}{z'}\right) = \arg z-\arg z' \ [2\pi],
  \qquad \arg\!\left(\overline z\right) = -\arg z \ [2\pi].
\]
\end{propriete}

% Figure adaptée des constructions du dossier TCD/Complexe/Lesona.tex.
\begin{visualisation}[multiplier, c'est tourner et dilater]
\begin{center}
\begin{tikzpicture}[x=1.05cm,y=1.05cm,font=\footnotesize]
  \coordinate (O) at (0,0);
  \coordinate (Z) at (25:2.15);
  \coordinate (U) at (70:1.55);
  \coordinate (P) at (95:{2.15*1.55});
  \draw[bmtGris,-{Stealth[length=2.2mm]}] (-.25,0)--(3.35,0);
  \draw[bmtGris,-{Stealth[length=2.2mm]}] (0,-.25)--(0,3.65);
  \draw[bmtIndigo,line width=1.05pt,-{Stealth[length=2.3mm]}] (O)--(Z)
    node[pos=.72,below right] {$z$};
  \draw[bmtVetiver,line width=1.05pt,-{Stealth[length=2.3mm]}] (O)--(U)
    node[pos=.85,right=5pt] {$z'$};
  \draw[bmtLaterite,line width=1.25pt,-{Stealth[length=2.5mm]}] (O)--(P)
    node[pos=.78,left] {$zz'$};
  \draw[bmtIndigo] (0:.68) arc (0:25:.68);
  \node[bmtIndigo] at (12.5:.92) {$\alpha$};
  \draw[bmtVetiver] (25:1.02) arc (25:95:1.02);
  \node[bmtVetiver] at (45:1.27) {$\beta$};
  \draw[bmtGris,dashed] (Z) arc (25:95:2.15);
  \fill[bmtEncre] (O) circle (1.7pt) node[below left] {$O$};
\end{tikzpicture}
\end{center}

Si $z=re^{i\alpha}$ et $z'=r'e^{i\beta}$, le produit a pour module
$rr'$ et pour argument $\alpha+\beta$. Sur la figure, le passage de $z$ à
$zz'$ combine donc une rotation d'angle $\beta$ et une multiplication des
longueurs par $r'$. Cette lecture géométrique est celle que l'on réutilisera
dans le chapitre sur les similitudes.
\end{visualisation}

\begin{propriete}[lecture géométrique d'un quotient]
Pour trois points distincts $A$, $B$, $C$ :
\[
  \left|\frac{z_{C}-z_{A}}{z_{B}-z_{A}}\right| = \frac{AC}{AB},
  \qquad
  \arg\!\left(\frac{z_{C}-z_{A}}{z_{B}-z_{A}}\right)
  = \left(\vect{AB},\vect{AC}\right) \ [2\pi].
\]
\end{propriete}

\begin{methode}{reconnaître la nature d'un triangle}
On calcule le quotient $\dfrac{z_{C}-z_{A}}{z_{B}-z_{A}}$ et l'on interprète.
Module $1$ : le triangle est isocèle en $A$. Argument $\frac{\pi}{2}$ : il est
rectangle en $A$. Quotient égal à $i$ : rectangle et isocèle en $A$. Quotient
réel : les trois points sont alignés.
\end{methode}

\begin{exemple}[nature d'un triangle]
Soient $A(2i)$, $B(3+i)$ et $C(-1-i)$. Alors
\[
  \frac{z_{B}-z_{A}}{z_{C}-z_{A}} = \frac{3-i}{-1-3i}
  = \frac{(3-i)(-1+3i)}{1+9} = \frac{-3+9i+i+3}{10} = i .
\]
Le module vaut $1$ et l'argument $\frac{\pi}{2}$ : le triangle $ABC$ est
rectangle et isocèle en $A$.
\end{exemple}

\begin{entrainement}[module et argument]
\exo[\niveau{1}] Donner module et argument de $1+i$, de $-\sqrt3+i$ et de
$-2i$.

\exo[\niveau{1}] Écrire $2\left(\cos\frac{\pi}{6}+i\sin\frac{\pi}{6}\right)$
sous forme algébrique.

\exo[\niveau{2}] Déterminer l'ensemble des points $M$ d'affixe $z$ tels que
$\abs{z-1+2i} = 3$.

\exo[\niveau{2}] Déterminer l'ensemble des points $M$ tels que
$\abs{z-1} = \abs{z+i}$.

\exo[\niveau{3}] Soient $A(1)$, $B(i)$ et $C(-1)$. Déterminer la nature du
triangle $ABC$ par le calcul d'un quotient.
\end{entrainement}

\begin{corrige}
\textbf{6.} $\abs{1+i} = \sqrt2$ d'argument $\frac{\pi}{4}$ ;
$\abs{-\sqrt3+i} = 2$ d'argument $\frac{5\pi}{6}$ ; $\abs{-2i} = 2$
d'argument $-\frac{\pi}{2}$.

\textbf{7.} $2\left(\cos\frac{\pi}{6}+i\sin\frac{\pi}{6}\right)
= 2\left(\frac{\sqrt3}{2}+\frac12 i\right) = \sqrt3+i$.

\textbf{8.} L'égalité s'écrit $\abs{z-(1-2i)} = 3$ : c'est le cercle de
centre le point d'affixe $1-2i$ et de rayon $3$.

\textbf{9.} Avec $z = x+iy$, l'égalité devient
$(x-1)^{2}+y^{2} = x^{2}+(y+1)^{2}$, soit $-2x = 2y$, c'est-à-dire $y = -x$.
C'est la médiatrice du segment joignant les points d'affixes $1$ et $-i$.

\textbf{10.} $\dfrac{z_{C}-z_{A}}{z_{B}-z_{A}} = \dfrac{-2}{i-1} = 1+i$, de
module $\sqrt2$ et d'argument $\frac{\pi}{4}$. Donc $AC = \sqrt2\,AB$ et
l'angle en $A$ vaut $\frac{\pi}{4}$. Comme $AB = \abs{i-1} = \sqrt2$,
$BC = \abs{-1-i} = \sqrt2$ et $AC = 2$, on a $AB = BC$ et
$AB^{2}+BC^{2} = AC^{2}$ : le triangle est rectangle isocèle en $B$.
\end{corrige}

% ===========================================================================
\section{Formes exponentielle et polaire}

\begin{definition}[notation exponentielle]
Pour tout réel $\theta$, on pose
$e^{i\theta} = \cos\theta+i\sin\theta$. Tout complexe non nul s'écrit alors
$z = re^{i\theta}$ avec $r = \abs z>0$ : c'est la \textbf{forme
exponentielle}. La notation $[r\,;\theta]$ est la \textbf{forme polaire}.
\end{definition}

\begin{propriete}[pourquoi cette notation]
\[
  e^{i\theta}\times e^{i\theta'} = e^{i(\theta+\theta')}, \qquad
  \frac{1}{e^{i\theta}} = e^{-i\theta}, \qquad
  \left(e^{i\theta}\right)^{n} = e^{in\theta}.
\]
La notation exponentielle est donc cohérente avec les règles de calcul de
l'exponentielle réelle, ce qui justifie son emploi.
\end{propriete}

\begin{demonstration}
Le produit se déduit des formules d'addition :
\[
  (\cos\theta+i\sin\theta)(\cos\theta'+i\sin\theta')
  = \left(\cos\theta\cos\theta'-\sin\theta\sin\theta'\right)
  +i\left(\sin\theta\cos\theta'+\cos\theta\sin\theta'\right),
\]
c'est-à-dire $\cos(\theta+\theta')+i\sin(\theta+\theta')$.
\end{demonstration}

% ===========================================================================
\section{Équation du second degré dans l'ensemble des complexes}

\begin{theoreme}[racines carrées d'un réel négatif]
Si $a>0$, les solutions de $z^{2} = -a$ sont $i\sqrt a$ et $-i\sqrt a$.
\end{theoreme}

\begin{theoreme}[équation du second degré à coefficients réels]
Soit $az^{2}+bz+c = 0$ avec $a$, $b$, $c$ réels et $a \neq 0$, de
discriminant $\Delta = b^{2}-4ac$.
\begin{itemize}[leftmargin=*, itemsep=0.1em]
  \item Si $\Delta>0$, deux solutions réelles
        $\dfrac{-b\pm\sqrt\Delta}{2a}$.
  \item Si $\Delta = 0$, une solution double $-\dfrac{b}{2a}$.
  \item Si $\Delta<0$, deux solutions complexes conjuguées
        $\dfrac{-b\pm i\sqrt{-\Delta}}{2a}$.
\end{itemize}
\end{theoreme}

\begin{propriete}[coefficients complexes]
Lorsque les coefficients sont complexes, la formule reste valable à condition
de remplacer $\sqrt\Delta$ par \emph{une} racine carrée $\delta$ de $\Delta$,
c'est-à-dire un complexe vérifiant $\delta^{2} = \Delta$. Les solutions sont
alors $\dfrac{-b\pm\delta}{2a}$.
\end{propriete}

\begin{methode}{extraire une racine carrée d'un complexe}
Pour trouver $\delta = a+ib$ tel que $\delta^{2} = \Delta = p+iq$, on écrit
trois équations : $a^{2}-b^{2} = p$, $2ab = q$, et
$a^{2}+b^{2} = \abs\Delta$. Les deux dernières donnent $a^{2}$ et $b^{2}$ par
somme et différence, la deuxième fixe le signe relatif de $a$ et $b$.
\end{methode}

\begin{exemple}[une équation à coefficients complexes]
Résolvons $z^{2}-(5-3i)z+4-7i = 0$.

$\Delta = (5-3i)^{2}-4(4-7i) = (16-30i)-16+28i = -2i$.

Cherchons $\delta = a+ib$ avec $\delta^{2} = -2i$ : $a^{2}-b^{2} = 0$,
$2ab = -2$ et $a^{2}+b^{2} = 2$. On obtient $a = 1$, $b = -1$, soit
$\delta = 1-i$ — on vérifie $(1-i)^{2} = -2i$.

Les solutions sont
\[
  z_{1} = \frac{(5-3i)+(1-i)}{2} = 3-2i,
  \qquad
  z_{2} = \frac{(5-3i)-(1-i)}{2} = 2-i .
\]
\end{exemple}

\begin{entrainement}[équations]
\exo[\niveau{1}] Résoudre dans $\C$ : $z^{2}+4 = 0$, puis $z^{2}-2z+5 = 0$.

\exo[\niveau{2}] Résoudre $z^{2}-(3+i)z+4+3i = 0$.

\exo[\niveau{2}] Déterminer les racines carrées de $3+4i$.

\exo[\niveau{3}] Résoudre $z^{2}-2z\cos\theta+1 = 0$, $\theta$ réel, et donner
les solutions sous forme exponentielle.
\end{entrainement}

\begin{corrige}
\textbf{11.} $z^{2}+4 = 0$ donne $z = 2i$ ou $z = -2i$. Pour
$z^{2}-2z+5 = 0$, le discriminant vaut $-16$, de racines carrées $\pm4i$,
d'où $z = 1+2i$ ou $z = 1-2i$.

\textbf{12.} $\Delta = (3+i)^{2}-4(4+3i) = 8+6i-16-12i = -8-6i$. On cherche
$\delta = a+ib$ tel que $\delta^{2} = \Delta$ : le système
$a^{2}-b^{2} = -8$, $2ab = -6$, $a^{2}+b^{2} = \abs\Delta = 10$ donne
$a^{2} = 1$ et $b^{2} = 9$ avec $ab<0$, soit $\delta = 1-3i$. Les solutions
sont alors $\dfrac{(3+i)\pm(1-3i)}{2}$, c'est-à-dire $2-i$ et $1+2i$. On
vérifie que leur somme vaut $3+i$ et leur produit $4+3i$.

\textbf{13.} Si $(a+ib)^{2} = 3+4i$, alors $a^{2}-b^{2} = 3$, $2ab = 4$ et
$a^{2}+b^{2} = 5$. D'où $a^{2} = 4$, $b^{2} = 1$, avec $ab>0$ : les racines
carrées sont $2+i$ et $-2-i$.

\textbf{14.} $\Delta = 4\cos^{2}\theta-4 = -4\sin^{2}\theta$, de racines
carrées $\pm2i\sin\theta$. Les solutions sont donc
$z = \cos\theta \pm i\sin\theta$, c'est-à-dire $e^{i\theta}$ et
$e^{-i\theta}$ : deux nombres de module $1$, conjugués l'un de l'autre.
\end{corrige}

% ===========================================================================
\section{Résolution algébrique et factorisation de polynômes}

\begin{propriete}[factorisation par une racine]
Si $z_{0}$ est racine du polynôme $P$ de degré $n$, alors $P(z)$ se factorise
sous la forme $(z-z_{0})\,Q(z)$, où $Q$ est de degré $n-1$.
\end{propriete}

\begin{methode}{résoudre une équation du troisième degré}
L'énoncé fournit toujours une racine, ou demande de la vérifier. On divise
alors $P$ par $z-z_{0}$ — par la division euclidienne posée, ou par
identification des coefficients —, et l'on résout l'équation du second degré
obtenue.
\end{methode}

\begin{exemple}[une équation du troisième degré]
Résolvons $P(z) = z^{3}+(3-3i)z^{2}+(1-6i)z-1-3i = 0$.

\emph{Une racine évidente.}
$P(-1) = -1+(3-3i)-(1-6i)-1-3i = 0$ : $-1$ est racine.

\emph{Factorisation.} La division donne
$P(z) = (z+1)\left[z^{2}+(2-3i)z-1-3i\right]$.

\emph{Second degré.}
$\Delta = (2-3i)^{2}-4(-1-3i) = (-5-12i)+4+12i = -1$, dont $i$ est une racine
carrée. Les solutions sont
\[
  \frac{-(2-3i)+i}{2} = -1+2i
  \qquad\text{et}\qquad
  \frac{-(2-3i)-i}{2} = -1+i .
\]
L'ensemble des solutions est $\{-1\,;\,-1+i\,;\,-1+2i\}$.
\end{exemple}

% ===========================================================================
\section{Formules de Moivre et d'Euler}

\begin{theoreme}[formule de Moivre]
Pour tout réel $\theta$ et tout entier $n$,
\[
  \left(\cos\theta+i\sin\theta\right)^{n} = \cos(n\theta)+i\sin(n\theta).
\]
\end{theoreme}

\begin{demonstration}
Récurrence sur $n$. Vrai pour $n = 0$. Si la formule vaut au rang $n$, alors
en multipliant par $\cos\theta+i\sin\theta$ et en utilisant les formules
d'addition, on obtient $\cos((n+1)\theta)+i\sin((n+1)\theta)$.
\end{demonstration}

\begin{theoreme}[formules d'Euler]
\[
  \cos\theta = \frac{e^{i\theta}+e^{-i\theta}}{2},
  \qquad
  \sin\theta = \frac{e^{i\theta}-e^{-i\theta}}{2i}.
\]
\end{theoreme}

\begin{exemple}[linéariser $\cos^{3}\theta$]
\[
  \cos^{3}\theta
  = \left(\frac{e^{i\theta}+e^{-i\theta}}{2}\right)^{3}
  = \frac18\left(e^{3i\theta}+3e^{i\theta}+3e^{-i\theta}+e^{-3i\theta}\right)
  = \frac{\cos3\theta+3\cos\theta}{4}.
\]
La linéarisation transforme une puissance en somme : c'est exactement ce
qu'il faut pour calculer une primitive de $\cos^{3}$.
\end{exemple}

\begin{entrainement}[Moivre et Euler]
\exo[\niveau{1}] Calculer $(1+i)^{8}$ en passant par la forme exponentielle.

\exo[\niveau{2}] Exprimer $\cos3\theta$ en fonction de $\cos\theta$.

\exo[\niveau{2}] Linéariser $\sin^{2}\theta$ puis $\sin^{3}\theta$.

\exo[\niveau{3}] En déduire une primitive de $\theta \mapsto \sin^{3}\theta$.
\end{entrainement}

\begin{corrige}
\textbf{15.} $1+i = \sqrt2\,e^{i\pi/4}$, donc
$(1+i)^{8} = \left(\sqrt2\right)^{8}e^{2i\pi} = 16$.

\textbf{16.} La formule de Moivre donne
$\cos3\theta+i\sin3\theta = \left(\cos\theta+i\sin\theta\right)^{3}$. En
développant et en identifiant les parties réelles,
$\cos3\theta = \cos^{3}\theta-3\cos\theta\sin^{2}\theta$, puis, en remplaçant
$\sin^{2}\theta$ par $1-\cos^{2}\theta$,
$\cos3\theta = 4\cos^{3}\theta-3\cos\theta$.

\textbf{17.} Avec les formules d'Euler,
$\sin^{2}\theta = \dfrac{1-\cos2\theta}{2}$ et
$\sin^{3}\theta = \dfrac{3\sin\theta-\sin3\theta}{4}$.

\textbf{18.} La linéarisation transforme la primitive en une somme de deux
primitives immédiates :
\[
  \int \sin^{3}\theta \,\mathrm{d}\theta
  = \frac14\int\left(3\sin\theta-\sin3\theta\right)\,\mathrm{d}\theta
  = -\frac34\cos\theta+\frac{1}{12}\cos3\theta+k .
\]
\end{corrige}

% ===========================================================================
\section{Racines n-ièmes d'un nombre complexe non nul}

\begin{theoreme}[racines n-ièmes]
Soit $Z = Re^{i\varphi}$ un complexe non nul et $n \geq 1$. L'équation
$z^{n} = Z$ admet exactement $n$ solutions :
\[
  z_{k} = R^{1/n}\,e^{i\left(\frac{\varphi}{n}+\frac{2k\pi}{n}\right)},
  \qquad k = 0,1,\dots,n-1 .
\]
Leurs images sont les sommets d'un polygone régulier à $n$ côtés inscrit dans
le cercle de centre $O$ et de rayon $R^{1/n}$.
\end{theoreme}

% Figure adaptée des polygones réguliers de TCD/Appendix/Polygones.tex.
\begin{visualisation}[les racines de l'unité sur un cercle]
\begin{center}
\begin{tikzpicture}[scale=1.02,font=\footnotesize]
  \draw[bmtGris!55] (0,0) circle (2);
  \draw[bmtGris,-{Stealth[length=2.1mm]}] (-2.35,0)--(2.55,0);
  \draw[bmtGris,-{Stealth[length=2.1mm]}] (0,-2.35)--(0,2.5);
  \draw[bmtAccent,line width=1.05pt]
    (0:2)--(60:2)--(120:2)--(180:2)--(240:2)--(300:2)--cycle;
  \foreach \k in {0,...,5}{
    \fill[bmtLaterite] ({60*\k}:2) circle (2.2pt);
    \node[fill=bmtPapier,inner sep=.6pt] at ({60*\k}:2.32) {$\omega^{\k}$};
  }
  \fill[bmtEncre] (0,0) circle (1.6pt) node[below left] {$O$};
  \draw[bmtVetiver,-{Stealth[length=2mm]}] (0:.72) arc (0:60:.72);
  \node[bmtVetiver] at (30:1.02) {$\frac{2\pi}{6}$};
\end{tikzpicture}
\end{center}

Pour $z^{6}=1$, les six solutions ont le même module $1$ et deux arguments
consécutifs diffèrent de $\frac{2\pi}{6}$. Elles forment donc un hexagone
régulier. La même construction fonctionne pour tout $n$ : on partage le tour
complet en $n$ angles égaux.
\end{visualisation}

\begin{demonstration}
En écrivant $z = re^{i\theta}$, l'équation devient $r^{n}e^{in\theta}
= Re^{i\varphi}$, ce qui équivaut à $r^{n} = R$ et
$n\theta = \varphi+2k\pi$. La première égalité donne $r = R^{1/n}$, la seconde
$\theta = \frac{\varphi}{n}+\frac{2k\pi}{n}$ ; les valeurs $k = 0,\dots,n-1$
donnent $n$ arguments distincts modulo $2\pi$, les autres les répètent.
\end{demonstration}

\begin{entrainement}[racines n-ièmes]
\exo[\niveau{1}] Résoudre $z^{3} = 1$ et représenter les solutions.

\exo[\niveau{2}] Résoudre $z^{4} = -16$.

\exo[\niveau{2}] Résoudre $z^{3} = 8i$.

\exo[\niveau{3}] Montrer que la somme des racines $n$-ièmes de l'unité est
nulle pour $n \geq 2$.
\end{entrainement}

\begin{corrige}
\textbf{19.} Les solutions de $z^{3} = 1$ sont $1$, $e^{2i\pi/3}$ et
$e^{-2i\pi/3}$ : trois points régulièrement répartis sur le cercle unité, aux
sommets d'un triangle équilatéral.

\textbf{20.} $-16 = 16\,e^{i\pi}$, donc $z = 2\,e^{i(\pi+2k\pi)/4}$ pour
$k \in \{0,1,2,3\}$ : les solutions sont $\sqrt2+i\sqrt2$,
$-\sqrt2+i\sqrt2$, $-\sqrt2-i\sqrt2$ et $\sqrt2-i\sqrt2$.

\textbf{21.} $8i = 8\,e^{i\pi/2}$, donc
$z = 2\,e^{i\left(\frac{\pi}{6}+\frac{2k\pi}{3}\right)}$ : les solutions sont
$\sqrt3+i$, $-\sqrt3+i$ et $-2i$.

\textbf{22.} Les racines $n$-ièmes de l'unité sont les $\omega^{k}$ pour
$k$ de $0$ à $n-1$, où $\omega = e^{2i\pi/n} \neq 1$. Leur somme est celle
d'une suite géométrique de raison $\omega$ :
$\dfrac{\omega^{n}-1}{\omega-1} = \dfrac{1-1}{\omega-1} = 0$.
\end{corrige}

% ===========================================================================
\begin{annales}
Les nombres complexes ouvrent presque toujours par une équation, puis
basculent vers la géométrie. Les énoncés qui suivent sont repris de sujets
réellement posés ; leur provenance est indiquée à chaque fois.

\exoann{Baccalauréat, Madagascar, série S, session 2021 — problème 1, partie C}
Dans $\C$, on considère le polynôme
\[
  P(z) = z^{3}+(3-3i)z^{2}+(1-6i)z-1-3i .
\]
Calculer $P(-1)$ et en déduire toutes les solutions dans $\C$ de l'équation
$P(z) = 0$.

\exoann{Baccalauréat, Mauritanie, série C, session 2022 — exercice 2}
On pose $P(z) = z^{3}-(5+6i)z^{2}+(-6+16i)z+20+10i$.
\begin{enumerate}[label=\textbf{\arabic*)}, leftmargin=*, itemsep=0.15em]
  \item Calculer $P(-1+i)$ et factoriser $P(z)$ sous la forme
        $(z+1-i)\left(z^{2}+az+b\right)$.
  \item Résoudre $P(z) = 0$. On note $z_{1}$, $z_{2}$, $z_{3}$ les solutions,
        rangées de sorte que $\abs{z_{1}}<\abs{z_{2}}<\abs{z_{3}}$.
\end{enumerate}

\exoann{Baccalauréat, Burkina Faso, série D, session 2023 — exercice 1}
Soit $P(Z) = Z^{3}-(2+2i)Z^{2}-2Z-8+4i$.
\begin{enumerate}[label=\textbf{\arabic*)}, leftmargin=*, itemsep=0.15em]
  \item Vérifier que $2i$ est racine de $P(Z)$.
  \item Factoriser $P(Z)$, puis résoudre l'équation $P(Z) = 0$.
\end{enumerate}

\exoann{Baccalauréat, Burkina Faso, série D, session 2023 — exercice 1, suite}
Dans le plan complexe, on place $A(2i)$, $B(3+i)$ et $C(-1-i)$. Calculer
$\dfrac{z_{B}-z_{A}}{z_{C}-z_{A}}$ et en déduire la nature du triangle $ABC$.

\exoann{Baccalauréat, Cameroun, séries D et TI, session 2024 — exercice 1}
Résoudre dans $\C$ l'équation
$z^{2}-(5-3i)z+4-7i = 0$.

\exoann{Baccalauréat, Cameroun, série C, session 2023 — exercice 2}
Soit $\alpha$ un réel et $(E)$ l'équation d'inconnue complexe
\[
  z^{2}+\bigl[-3\cos\alpha-1+i(3-5\sin\alpha)\bigr]z
  +5\sin\alpha-2+i(-3\cos\alpha-1) = 0 .
\]
\begin{enumerate}[label=\textbf{\arabic*)}, leftmargin=*, itemsep=0.15em]
  \item Montrer que $-i$ est solution de $(E)$.
  \item Déterminer l'autre solution $z_{2}$.
\end{enumerate}

\exoann{Baccalauréat, Cameroun, série C, session 2023 — exercice 2, suite}
Déterminer l'ensemble des points $A_{\alpha}$ d'affixe
$z_{\alpha} = 3\cos\alpha+1+i\left(-2+5\sin\alpha\right)$ lorsque $\alpha$
décrit $\R$.

\exoann{Concours d'entrée, ENI Fianarantsoa, session 2008-2009 — exercice 1}
Soit l'équation d'inconnue complexe $z$ :
\[
  z^{2}-(\alpha+3i+4)z+2\alpha i-1 = 0, \qquad \alpha \in \C .
\]
\begin{enumerate}[label=\textbf{\alph*)}, leftmargin=*, itemsep=0.15em]
  \item Déterminer $\alpha$ pour que l'équation admette deux racines
        complexes conjuguées, puis calculer ces racines.
  \item Dans un autre cas, si l'une des racines est $i$, calculer $\alpha$ et
        l'autre racine.
\end{enumerate}

\exoann{Baccalauréat blanc, lycée Philibert Tsiranana, Terminale S, 2021-2022 — problème 1, partie C}
Dans le plan orienté, $OAB$ est un triangle tel que $OA = OB = a$ avec $a>0$
et $\left(\vect{OA},\vect{OB}\right) = \frac{\pi}{2}$. On note $I$ le milieu
de $[AB]$ et $G$ le point défini par
$2\,\vect{OG} = -\vect{OA}-\vect{OB}$. Le plan est rapporté au repère
complexe $\left(O\,;\vect{OA},\vect{OB}\right)$.

Déterminer les affixes des points $A$, $B$, $I$ et $G$.
\end{annales}

\begin{corrige}[annales]
\textbf{1.} $P(-1) = -1+(3-3i)-(1-6i)-1-3i = 0$. La division par $z+1$ donne
$P(z) = (z+1)\left[z^{2}+(2-3i)z-1-3i\right]$. Le discriminant vaut
$(2-3i)^{2}+4+12i = -1$, dont $i$ est une racine carrée. Les solutions sont
$-1$, $-1+i$ et $-1+2i$.

\textbf{2.} Avec $z = -1+i$ : $z^{2} = -2i$ et $z^{3} = 2+2i$, d'où
$P(-1+i) = (2+2i)+(12-10i)+(-10-22i)+(20+10i) = 0$. La division donne
$P(z) = (z+1-i)\left[z^{2}-(6+5i)z+5+15i\right]$. Le discriminant vaut
$(6+5i)^{2}-4(5+15i) = -9$, de racine carrée $3i$, d'où les solutions
$3+4i$ et $3+i$. Enfin $\abs{-1+i} = \sqrt2$, $\abs{3+i} = \sqrt{10}$ et
$\abs{3+4i} = 5$ : $z_{1} = -1+i$, $z_{2} = 3+i$, $z_{3} = 3+4i$.

\textbf{3.} $(2i)^{3} = -8i$ et $(2i)^{2} = -4$, donc
$P(2i) = -8i+(8+8i)-4i-8+4i = 0$. La division donne
$P(Z) = (Z-2i)\left(Z^{2}-2Z-2-4i\right)$. Le discriminant vaut $12+16i$, de
racine carrée $4+2i$ — car $(4+2i)^{2} = 12+16i$ —, d'où les solutions
$3+i$ et $-1-i$. L'ensemble des solutions est
$\{2i\,;\,3+i\,;\,-1-i\}$.

\textbf{4.} $\dfrac{3-i}{-1-3i} = \dfrac{(3-i)(-1+3i)}{10} = \dfrac{10i}{10}
= i$. Le module vaut $1$ et l'argument $\frac{\pi}{2}$ : $AB = AC$ et
$\left(\vect{AC},\vect{AB}\right) = \frac{\pi}{2}$. Le triangle est rectangle
et isocèle en $A$.

\textbf{5.} $\Delta = (5-3i)^{2}-4(4-7i) = -2i$, dont $1-i$ est une racine
carrée. Les solutions sont $3-2i$ et $2-i$.

\textbf{6.} Pour $z = -i$ : $z^{2} = -1$, le deuxième terme vaut
$i(3\cos\alpha+1)+(3-5\sin\alpha)$ et le troisième
$5\sin\alpha-2+i(-3\cos\alpha-1)$. La somme des parties réelles est
$-1+3-5\sin\alpha+5\sin\alpha-2 = 0$, celle des parties imaginaires
$(3\cos\alpha+1)-(3\cos\alpha+1) = 0$ : $-i$ est bien solution. Le produit
des racines vaut $5\sin\alpha-2+i(-3\cos\alpha-1)$, donc
\[
  z_{2} = \frac{5\sin\alpha-2+i(-3\cos\alpha-1)}{-i}
  = 3\cos\alpha+1+i\left(5\sin\alpha-2\right).
\]

\textbf{7.} En posant $x = 3\cos\alpha+1$ et $y = 5\sin\alpha-2$, on obtient
$\left(\frac{x-1}{3}\right)^{2}+\left(\frac{y+2}{5}\right)^{2} = 1$ :
l'ensemble décrit est l'ellipse de centre $\Omega(1\,;-2)$, de demi-axes $3$
et $5$.

\textbf{8. a)} Deux racines conjuguées imposent une somme et un produit
réels. En posant $\alpha = a+ib$ : la somme $\alpha+3i+4 = (a+4)+(b+3)i$ est
réelle si $b = -3$ ; le produit $2\alpha i-1 = -2b-1+2ai$ est réel si
$a = 0$. Donc $\alpha = -3i$, l'équation devient $z^{2}-4z+5 = 0$, de
solutions $2+i$ et $2-i$.
\textbf{b)} Si $i$ est racine, $-1-(\alpha+3i+4)i+2\alpha i-1 = 0$ donne
$1+(\alpha-4)i = 0$, soit $\alpha = 4+i$. La somme des racines vaut alors
$\alpha+3i+4 = 8+4i$, et l'autre racine est $8+4i-i = 8+3i$.

\textbf{9.} Dans le repère $\left(O\,;\vect{OA},\vect{OB}\right)$, les
vecteurs de base sont $\vect{OA}$ et $\vect{OB}$ : donc $z_{A} = 1$ et
$z_{B} = i$. Le milieu de $[AB]$ a pour affixe $\frac{1+i}{2}$. Enfin la
relation $2\,\vect{OG} = -\vect{OA}-\vect{OB}$ se traduit sur les affixes par
$2\left(z_{G}-0\right) = -(1-0)-(i-0)$, d'où
\[
  z_{G} = -\frac{1+i}{2}.
\]
\end{corrige}

% ===========================================================================
\begin{jecherche}[le pavage d'une place]
Une place circulaire de rayon $8$ mètres doit recevoir six lampadaires
disposés aux sommets d'un hexagone régulier inscrit dans un cercle de rayon
$6$ mètres, centré au centre de la place. Le plan est rapporté à un repère
orthonormé direct d'origine le centre de la place, l'unité valant un mètre.
Le premier lampadaire est placé au point d'affixe $6$.

\begin{enumerate}[label=\textbf{\arabic*.}, leftmargin=*, itemsep=0.3em]
  \item Justifier que les affixes des six lampadaires sont les solutions de
        l'équation $z^{6} = 6^{6}$.
  \item Déterminer ces six affixes sous forme exponentielle, puis sous forme
        algébrique.
  \item Calculer la distance entre deux lampadaires voisins.
  \item Le service technique veut ajouter un septième lampadaire au centre.
        Justifier que la somme des affixes des six premiers est nulle, et
        interpréter.
  \item On fait tourner l'ensemble du dispositif d'un angle de
        $\frac{\pi}{6}$ autour du centre. Quelle est la nouvelle affixe du
        premier lampadaire ?
  \item Un câble relie chaque lampadaire au centre. Quelle longueur totale de
        câble faut-il ?
\end{enumerate}
\end{jecherche}

\begin{corrige}[le pavage d'une place]
\textbf{1.} Les six points sont sur le cercle de rayon $6$ et se déduisent
les uns des autres par des rotations d'angle $\frac{2\pi}{6}$ : leurs affixes
sont donc de module $6$ et d'arguments $\frac{2k\pi}{6}$, ce qui caractérise
exactement les racines sixièmes de $6^{6}$.

\textbf{2.} $z_{k} = 6e^{ik\pi/3}$ pour $k = 0,\dots,5$, soit
\[
  6, \quad 3+3i\sqrt3, \quad -3+3i\sqrt3, \quad -6, \quad -3-3i\sqrt3,
  \quad 3-3i\sqrt3 .
\]

\textbf{3.} $\abs{z_{1}-z_{0}} = \abs{6e^{i\pi/3}-6} = 6$ : dans un hexagone
régulier, le côté est égal au rayon. Les lampadaires voisins sont donc à six
mètres l'un de l'autre.

\textbf{4.} La somme des racines $n$-ièmes vaut $0$ dès que $n \geq 2$ :
c'est la somme d'une suite géométrique de raison $e^{i\pi/3} \neq 1$, égale à
$6\dfrac{1-e^{2i\pi}}{1-e^{i\pi/3}} = 0$. Géométriquement, le centre de
gravité des six lampadaires est le centre de la place : le septième y sera
donc parfaitement centré.

\textbf{5.} Une rotation d'angle $\frac{\pi}{6}$ multiplie l'affixe par
$e^{i\pi/6}$ : le premier lampadaire passe en $6e^{i\pi/6}
= 3\sqrt3+3i$.

\textbf{6.} Chaque câble mesure $6$ mètres, soit $36$ mètres au total —
la rotation n'y change rien, puisqu'elle conserve les distances.
\end{corrige}

% ===========================================================================
\begin{jemeteste}
Répondre par vrai ou faux, en justifiant en une ligne.

\begin{enumerate}[label=\textbf{\arabic*.}, leftmargin=*, itemsep=0.28em]
  \item La partie imaginaire de $2-7i$ est $-7i$.
  \item $\abs{z+z'} = \abs z+\abs{z'}$ pour tous complexes.
  \item Si $z\overline z = 0$ alors $z = 0$.
  \item Une équation du second degré à coefficients complexes a toujours deux
        solutions, éventuellement confondues.
  \item Les solutions de $z^{4} = 1$ forment un carré.
  \item Si $\dfrac{z_{C}-z_{A}}{z_{B}-z_{A}}$ est réel, le triangle $ABC$ est
        rectangle en $A$.
\end{enumerate}
\end{jemeteste}

\begin{corrige}[je me teste]
\textbf{1.} Faux — elle vaut $-7$, qui est un réel.
\textbf{2.} Faux — c'est une inégalité : $\abs{z+z'} \leq \abs z+\abs{z'}$.
\textbf{3.} Vrai — $z\overline z = \abs z^{2}$.
\textbf{4.} Vrai — tout complexe admet des racines carrées.
\textbf{5.} Vrai — ce sont $1$, $i$, $-1$, $-i$, sommets d'un carré.
\textbf{6.} Faux — un quotient réel signifie que les points sont alignés ;
c'est un argument égal à $\frac{\pi}{2}$ qui donne l'angle droit.
\end{corrige}

% ===========================================================================
\begin{commeaubac}
\bmtsource{Baccalauréat, série S, session 2021 — problème 1, partie C}
Dans l'ensemble $\C$ des nombres complexes, on considère le polynôme $P$
défini par
\[
  P(z) = z^{3}+(3-3i)z^{2}+(1-6i)z-1-3i .
\]

\begin{enumerate}[label=\textbf{\arabic*.}, leftmargin=*, itemsep=0.25em]
  \item Calculer $P(-1)$. \hfill\textup{(0,25 pt)}
  \item En déduire toutes les solutions dans $\C$ de l'équation $P(z) = 0$.
        \hfill\textup{(0,5 pt)}
  \item Placer les images des trois solutions dans le plan complexe et
        déterminer la nature du triangle qu'elles forment.
        \hfill\textup{(0,5 pt)}
\end{enumerate}
\end{commeaubac}

\begin{corrige}[comme au baccalauréat]
\textbf{1.} $P(-1) = -1+3-3i-1+6i-1-3i = 0$.

\textbf{2.} La factorisation donne
$P(z) = (z+1)\left[z^{2}+(2-3i)z-1-3i\right]$, dont le discriminant vaut
$-1$. Les solutions sont $-1$, $-1+i$ et $-1+2i$.

\textbf{3.} Les trois images ont pour abscisse commune $-1$ et pour ordonnées
$0$, $1$ et $2$ : elles sont alignées sur la droite verticale d'équation
$x = -1$, régulièrement espacées. Il n'y a donc pas de triangle, et le point
$-1+i$ est le milieu des deux autres.
\end{corrige}

% ===========================================================================
\begin{concours}
\bmtsource{Concours d'entrée, ENI Fianarantsoa, session 2008-2009}
Soit l'équation d'inconnue complexe $z$ :
\[
  z^{2}-(\alpha+3i+4)z+2\alpha i-1 = 0, \qquad \alpha \in \C .
\]

\begin{enumerate}[label=\textbf{\arabic*.}, leftmargin=*, itemsep=0.25em]
  \item Déterminer le paramètre $\alpha$ pour que l'équation admette deux
        racines complexes conjuguées, puis calculer ces racines.
  \item Dans un autre cas, si l'une des racines est $i$, calculer $\alpha$ et
        l'autre racine.
\end{enumerate}
\end{concours}

\begin{corrige}[vers le concours]
\textbf{1.} Deux racines conjuguées $u$ et $\overline u$ ont pour somme
$2\Re(u)$ et pour produit $\abs u^{2}$ : toutes deux réelles. En écrivant
$\alpha = a+ib$, la somme $\alpha+3i+4 = (a+4)+(b+3)i$ est réelle si et
seulement si $b = -3$ ; le produit $2\alpha i-1 = (-2b-1)+2ai$ l'est si et
seulement si $a = 0$. Donc $\alpha = -3i$, l'équation devient
$z^{2}-4z+5 = 0$, de discriminant $-4$ : les racines sont $2+i$ et $2-i$.

\textbf{2.} Si $i$ est racine, alors
$i^{2}-(\alpha+3i+4)i+2\alpha i-1 = 0$, c'est-à-dire
$-1-\alpha i+3-4i+2\alpha i-1 = 0$, soit $1+(\alpha-4)i = 0$ et donc
$\alpha = 4+i$. La somme des racines vaut $\alpha+3i+4 = 8+4i$, d'où l'autre
racine $8+4i-i = 8+3i$. On vérifie le produit :
$i(8+3i) = -3+8i = 2\alpha i-1$.
\end{corrige}

% ===========================================================================
\begin{notepourlaclasse}
Le chapitre est long et les élèves l'abordent avec méfiance. Deux séances
suffisent pourtant à installer l'essentiel si l'on tient un principe : chaque
notion algébrique est immédiatement traduite dans le plan. Le module est une
distance, l'argument un angle, le conjugué une symétrie, le produit une
rotation-agrandissement. Un élève qui n'a pas fait ces quatre traductions ne
réussira pas le problème de géométrie du baccalauréat.

L'extraction d'une racine carrée d'un complexe est le point technique le plus
coûteux. Faites-la systématiser sous forme de trois équations, et exigez la
vérification finale : elle prend dix secondes et rattrape toutes les erreurs
de signe.

Enfin, les quotients $\frac{z_{C}-z_{A}}{z_{B}-z_{A}}$ méritent une séance à
eux seuls. C'est l'outil qui rend immédiate la reconnaissance des
configurations, et il sera repris tel quel dans les chapitres de
transformations. Ne le laissez pas passer comme une simple remarque.
\end{notepourlaclasse}

% =========================================================================
\begin{devoir}[2 heures]
Trois exercices indépendants, notés sur $20$. Le devoir est cumulatif : il porte sur ce chapitre et sur ceux qui le précèdent. Les énoncés sont repris de sessions réelles et assemblés ici en épreuve ; leur provenance est indiquée à chaque fois.

\exods{1}{Baccalauréat, Madagascar, série S, session 2026 — exercice 1, partie I}{5 points}
On considère les matrices
\[
  A = \mattrois{1}{0}{2}{0}{1}{3}{0}{0}{1}, \qquad
  B = \mattrois{2}{3}{5}{1}{0}{2}{0}{0}{-1}, \qquad I_{3}.
\]
\begin{enumerate}[label=\textbf{\arabic*)}, leftmargin=*, itemsep=0.15em]
  \item Calculer le déterminant de $A$ et conclure.
  \item Vérifier que $2A-A^{2} = I_{3}$, puis en déduire la matrice $P$ telle
        que $A \times P = I_{3}$.
\end{enumerate}

\exods{2}{Baccalauréat, Madagascar, série S, session 2022 — exercice 1, partie II}{7 points}
On considère la matrice
$A = \mattrois{0}{-1}{-1}{-1}{0}{-1}{-1}{-1}{0}$.
\begin{enumerate}[label=\textbf{\arabic*)}, leftmargin=*, itemsep=0.15em]
  \item Calculer $A^{2}$ et $A^{3}$.
  \item Vérifier que $4A+A^{2}-A^{3} = 4I_{3}$.
  \item En déduire que $A$ est inversible et donner $A^{-1}$.
\end{enumerate}

\exods{3}{Baccalauréat, Mauritanie, série C, session 2022 — exercice 2}{8 points}
On pose $P(z) = z^{3}-(5+6i)z^{2}+(-6+16i)z+20+10i$.
\begin{enumerate}[label=\textbf{\arabic*)}, leftmargin=*, itemsep=0.15em]
  \item Calculer $P(-1+i)$ et factoriser $P(z)$ sous la forme
        $(z+1-i)\left(z^{2}+az+b\right)$.
  \item Résoudre $P(z) = 0$. On note $z_{1}$, $z_{2}$, $z_{3}$ les solutions,
        rangées de sorte que $\abs{z_{1}}<\abs{z_{2}}<\abs{z_{3}}$.
\end{enumerate}

\end{devoir}

\begin{corrige}[devoir surveillé]
\textbf{Exercice 1.} La matrice $A$ est triangulaire supérieure : son déterminant est
le produit des coefficients diagonaux, soit $1$. Il est non nul, donc $A$ est
inversible. On calcule $A^{2} = \mattrois{1}{0}{4}{0}{1}{6}{0}{0}{1}$, d'où
$2A-A^{2} = I_{3}$. En factorisant, $A\left(2I_{3}-A\right) = I_{3}$ :
\[
  P = 2I_{3}-A = \mattrois{1}{0}{-2}{0}{1}{-3}{0}{0}{1}.
\]

\textbf{Exercice 2.} En posant $J$ la matrice dont tous les coefficients valent $1$, on
a $A = I_{3}-J$ et $J^{2} = 3J$. Donc
$A^{2} = I_{3}+J$ et $A^{3} = I_{3}-3J$. Il vient
$4A+A^{2}-A^{3} = 4I_{3}-4J+I_{3}+J-I_{3}+3J = 4I_{3}$. En factorisant,
$A\left(4I_{3}+A-A^{2}\right) = 4I_{3}$, donc $A$ est inversible et
\[
  A^{-1} = \tfrac14\left(4I_{3}+A-A^{2}\right) = I_{3}-\tfrac12 J
  = \tfrac12\mattrois{1}{-1}{-1}{-1}{1}{-1}{-1}{-1}{1}.
\]

\textbf{Exercice 3.} Avec $z = -1+i$ : $z^{2} = -2i$ et $z^{3} = 2+2i$, d'où
$P(-1+i) = (2+2i)+(12-10i)+(-10-22i)+(20+10i) = 0$. La division donne
$P(z) = (z+1-i)\left[z^{2}-(6+5i)z+5+15i\right]$. Le discriminant vaut
$(6+5i)^{2}-4(5+15i) = -9$, de racine carrée $3i$, d'où les solutions
$3+4i$ et $3+i$. Enfin $\abs{-1+i} = \sqrt2$, $\abs{3+i} = \sqrt{10}$ et
$\abs{3+4i} = 5$ : $z_{1} = -1+i$, $z_{2} = 3+i$, $z_{3} = 3+4i$.

\end{corrige}
